\chapter{Driver Assistance Portal and Geospatial Logic}

\section{Commuter Advisory Objectives}
The Driver Assistance Portal (\texttt{frontend/src/pages/DriverPortal.jsx} and \texttt{/api/driver\_advice}) provides actionable junction telemetry for approaching drivers. By arming drivers with instantaneous phase timers and corridor densities, commuters can anticipate braking patterns and avoid blocking yellow box junctions.

\section{Haversine Geodesic Distance Formulation}
The backend computes straight-line aerial distance using the spherical Haversine formula:
\begin{align}
\Delta\phi &= \phi_{\text{junction}} - \phi_{\text{user}} \quad (\text{radians}) \\
\Delta\lambda &= \lambda_{\text{junction}} - \lambda_{\text{user}} \quad (\text{radians}) \\
a &= \sin^2\left(\frac{\Delta\phi}{2}\right) + \cos(\phi_{\text{user}}) \cos(\phi_{\text{junction}}) \sin^2\left(\frac{\Delta\lambda}{2}\right) \\
c &= 2 \cdot \text{atan2}\left(\sqrt{a}, \sqrt{1 - a}\right) \\
d_{\text{km}} &= R_{\text{Earth}} \times c \quad \text{where } R_{\text{Earth}} = 6371.0\,\text{km}
\end{align}

\begin{lstlisting}[language=Python,caption={Haversine Distance Implementation in Python}]
def _calculate_haversine(lat1, lon1, lat2, lon2):
    R = 6371.0  # Earth radius in kilometers
    dlat = math.radians(lat2 - lat1)
    dlon = math.radians(lon2 - lon1)
    a = (math.sin(dlat / 2) ** 2 +
         math.cos(math.radians(lat1)) * math.cos(math.radians(lat2)) *
         math.sin(dlon / 2) ** 2)
    c = 2 * math.atan2(math.sqrt(a), math.sqrt(1 - a))
    return round(R * c, 2)
\end{lstlisting}

\section{Coordinate Validation and Boundary Enforcements}
In \texttt{backend/app.py} and \texttt{backend/app\_demo.py}, incoming parameters undergo strict sanity checks:
\begin{itemize}[leftmargin=*]
    \item String length is capped at 32 characters to prevent buffer overflow attacks or regex Denial of Service.
    \item Coordinates are converted to floating-point numbers with explicit guards against \texttt{NaN} and \texttt{Inf}.
    \item Geographic ranges are bounded to $-90.0 \le \text{lat} \le 90.0$ and $-180.0 \le \text{lon} \le 180.0$.
    \item If coordinates violate bounds, the API returns \texttt{distance\_telemetry.status: "COORDINATES\_OUT\_OF\_RANGE"} with \texttt{distance\_km: null}.
    \item If only latitude or only longitude is supplied, the API responds with \texttt{INCOMPLETE\_COORDINATES\_MISSING\_LON} or \texttt{INCOMPLETE\_COORDINATES\_MISSING\_LAT}.
\end{itemize}

\section{Straight-Line vs. Real Road Distance Disclosures}
\begin{conceptbox}[Geodesic Proximity vs. Road Driving Distance]
The Haversine calculation determines the theoretical great-circle aerial distance (as the crow flies) between commuter GPS coordinates and the junction center. It does \textbf{not} reflect:
\begin{enumerate}
    \item Actual street topology, elevated flyovers (e.g., the Electronic City Expressway), or mandatory U-turns.
    \item Physical driving distance, which is typically $20\%\text{--}45\%$ longer in complex urban networks.
    \item Dynamic queue delays.
\end{enumerate}
The Driver Assistance Portal renders an explicit notice informing commuters of these boundaries, preventing misleading travel time assumptions.
\end{conceptbox}
